\documentclass[11pt]{article}

\usepackage[margin=1in]{geometry}
\usepackage[T1]{fontenc}
\usepackage{lmodern}
\usepackage{microtype}
\usepackage{amsmath,amssymb}
\usepackage[dvipsnames]{xcolor}
\usepackage{enumitem}
\usepackage{needspace}
\usepackage[most]{tcolorbox}
\usepackage[hidelinks]{hyperref}

\hypersetup{
  pdftitle={Homework 2 Solutions, MATH 327, Fall 2026},
  pdfauthor={Manish Patnaik},
  pdfsubject={MATH 327 homework problems}
}

\definecolor{courseblue}{HTML}{1F4E79}
\definecolor{lightblue}{HTML}{EAF2F8}
\definecolor{boxborder}{HTML}{B8CBDC}
\definecolor{darkgray}{HTML}{333333}
\definecolor{lightyellow}{HTML}{FFF9DC}
\definecolor{yellowborder}{HTML}{E4D58C}

\setlength{\parindent}{0pt}
\setlength{\parskip}{0.55em}
\setlist[enumerate]{leftmargin=2em, itemsep=0.25em, topsep=0.3em}

\newtcolorbox{exercise}[1]{
  enhanced,
  colback=white,
  colframe=boxborder,
  colbacktitle=lightblue,
  coltitle=courseblue,
  fonttitle=\bfseries,
  title={Exercise #1},
  boxrule=0.6pt,
  arc=1mm,
  left=2.5mm,
  right=2.5mm,
  top=1.5mm,
  bottom=1.5mm,
  before skip=0.8em,
  after skip=0.8em
}

\newtcolorbox{solution}{
  enhanced,
  breakable,
  colback=lightyellow,
  colframe=yellowborder,
  colbacktitle=lightyellow,
  coltitle=darkgray,
  fonttitle=\bfseries,
  title={Solution},
  boxrule=0.6pt,
  arc=1mm,
  left=2.5mm,
  right=2.5mm,
  top=1.5mm,
  bottom=1.5mm,
  before skip=0.2em,
  after skip=1em
}

\newcommand{\chapterheading}[1]{%
  \par\vspace{0.5em}
  {\color{courseblue}\Large\bfseries #1}\par\vspace{0.35em}
  {\color{boxborder}\hrule height 0.6pt}\vspace{0.35em}
}

\newcommand{\sectionheading}[1]{%
  \Needspace{8\baselineskip}
  \vspace{0.25em}
  {\color{darkgray}\large\bfseries #1}\par
}

\begin{document}

\begin{center}
  \colorbox{lightblue}{%
    \begin{minipage}{0.92\textwidth}
      \vspace{0.6em}
      \begin{tabular*}{\textwidth}{@{\extracolsep{\fill}}lr}
        {\color{courseblue}\bfseries\LARGE Homework 2} &
        {\color{courseblue}\bfseries\LARGE MATH 327}\\[0.4em]
        \multicolumn{2}{l}{\color{darkgray}\bfseries Fall 2026}\\[0.35em]
        \multicolumn{2}{l}{\color{courseblue}\bfseries Questions and solutions}\\[0.35em]
        \multicolumn{2}{l}{\color{darkgray}Covers Lectures 8--13}\\[0.2em]
        \multicolumn{2}{l}{\color{darkgray}Suggested completion date: October 5}
      \end{tabular*}
      \vspace{0.5em}
    \end{minipage}}
\end{center}

\vspace{0.5em}

Unless otherwise indicated, the following exercises are from Michael Artin's
\emph{Algebra}, second edition.

% Solutions adapted from Homework (2025)/2025-Winter-MATH327-Homework2_solutions.tex.
% The index-3 example in 8.10 and the corrected M.12 argument were added
% with the instructor's permission.

\chapterheading{Chapter 2: Groups}

\sectionheading{Section 8: Cosets}

\begin{exercise}{8.1}
Let $H$ be the cyclic subgroup of the alternating group $A_4$ generated
by the permutation $(123)$. Exhibit the left and the right cosets of $H$
explicitly.
\end{exercise}

\begin{solution}
The group $A_4$ has $12$ elements and $H=\langle(123)\rangle$ has $3$,
so there are four left cosets and four right cosets. Direct computation gives
the left cosets
\begin{align*}
  H&=\{1,(123),(132)\},\\
  (12)(34)H&=\{(12)(34),(143),(243)\},\\
  (13)(24)H&=\{(13)(24),(142),(234)\},\\
  (14)(23)H&=\{(14)(23),(124),(134)\},
\end{align*}
and the right cosets
\begin{align*}
  H&=\{1,(123),(132)\},\\
  H(12)(34)&=\{(12)(34),(134),(234)\},\\
  H(13)(24)&=\{(13)(24),(124),(243)\},\\
  H(14)(23)&=\{(14)(23),(142),(143)\}.
\end{align*}
\end{solution}

\begin{exercise}{8.4}
Does a group of order $35$ contain an element of order $5$? Of order $7$?
\end{exercise}

\begin{solution}
If $G$ has an element $x$ of order $35$, then $x^7$ has order $5$ and
$x^5$ has order $7$. Otherwise, to show that $G$ has an element of order
$7$, suppose every nonidentity element has order $5$. Distinct subgroups
of order $5$ intersect only in the identity, so counting their elements
would give $35=1+4k$ for some integer $k$, a contradiction. The same
argument shows that $G$ has an element of order $5$: if every
nonidentity element had order $7$, then $35=1+6k$, also impossible.
\end{solution}

\begin{exercise}{8.10}
Prove that every subgroup of index $2$ is a normal subgroup, and show by
example that a subgroup of index $3$ need not be normal.
\end{exercise}

\begin{solution}
Suppose $[G:H]=2$, and let $a\notin H$. Since $aH\ne H$ and $Ha\ne H$,
the two-coset decompositions of $G$ give
\[
  G=H\sqcup aH=H\sqcup Ha.
\]
Thus $aH=Ha$ for every $a\notin H$; the same equality holds for $a\in H$.
Hence $H$ is normal in $G$.

For an index-$3$ example, take $G=S_3$ and $H=\{1,(12)\}$. Then
$[G:H]=3$, but $(123)(12)(132)=(23)\notin H$, so $H$ is not normal.
\end{solution}

\sectionheading{Section 9: Modular Arithmetic}

\begin{exercise}{9.6}
Prove the Chinese Remainder Theorem: Let $a,b,u,v$ be integers, and assume
that the greatest common divisor of $a$ and $b$ is $1$. Then there is an
integer $x$ such that
\[
  x\equiv u\pmod a
  \qquad\text{and}\qquad
  x\equiv v\pmod b.
\]
\textit{Hint.} Do the case $u=0$ and $v=1$ first.
\end{exercise}

\begin{solution}
Since $\gcd(a,b)=1$, choose integers $r,s$ with $ra+sb=1$. Set
$x=ra$ and $y=sb$. Then
\[
  x\equiv0\pmod a,\quad x\equiv1\pmod b,
  \qquad
  y\equiv1\pmod a,\quad y\equiv0\pmod b.
\]
Therefore $uy+vx$ satisfies both required congruences.
\end{solution}

\sectionheading{Miscellaneous Problems}

\begin{exercise}{M.12}
\textit{Postage stamp problem.} Let $a$ and $b$ be positive, relatively
prime integers.
\begin{enumerate}[label=\textbf{(\alph*)}]
  \item Prove that every sufficiently large positive integer $n$ can be
        obtained as $ra+sb$, where $r$ and $s$ are positive integers.
  \item Determine the largest integer that is not of this form.
\end{enumerate}
\end{exercise}

\begin{solution}
Since $\gcd(a,b)=1$, for any integer $n$ there is a unique
$r\in\{1,\ldots,b\}$ such that $ra\equiv n\pmod b$. Thus
\[
  s=\frac{n-ra}{b}
\]
is an integer. If $n>ab$, then $ra\leq ab<n$, so $s>0$; also $r>0$.
Hence every integer $n>ab$ has the required form $ra+sb$.

The integer $ab$ itself does not have this form with $r,s>0$. Indeed,
if $ab=ra+sb$, then reducing modulo $b$ gives $b\mid ra$. Since
$\gcd(a,b)=1$, we have $b\mid r$, so $r\geq b$ and $ra\geq ab$,
contradicting $sb>0$. Consequently, the largest integer not of the
required form is $\boxed{ab}$.
\end{solution}

\vfill
{\small\color{darkgray}Compiled \today.}

\end{document}
